\documentclass[11pt]{article}

\usepackage[margin=1in]{geometry}
\usepackage{amsmath,amssymb,amsthm}
\usepackage{mathtools}
\usepackage[colorlinks=true,linkcolor=black,citecolor=black,urlcolor=blue]{hyperref}
\usepackage{cleveref}
\usepackage{url}
\usepackage{enumitem}
\usepackage{booktabs}
\usepackage{array}
\usepackage{algorithm}
\usepackage{algpseudocode}

\newtheorem{theorem}{Theorem}[section]
\newtheorem{lemma}[theorem]{Lemma}
\newtheorem{proposition}[theorem]{Proposition}
\newtheorem{corollary}[theorem]{Corollary}
\newtheorem{definition}[theorem]{Definition}
\newtheorem{remark}[theorem]{Remark}
\newtheorem{observation}[theorem]{Observation}
\newtheorem{conjecture}[theorem]{Conjecture}

\DeclareMathOperator{\tr}{tr}
\DeclareMathOperator{\diag}{diag}
\DeclareMathOperator{\Id}{Id}
\DeclareMathOperator{\negl}{negl}
\newcommand{\norm}[1]{\lVert #1 \rVert}
\newcommand{\abs}[1]{\lvert #1 \rvert}
\newcommand{\frob}[1]{\norm{#1}_{\mathrm{F}}}
\newcommand{\opn}[1]{\norm{#1}_{\mathrm{op}}}
\newcommand{\dtv}{d_{\mathrm{TV}}}
\newcommand{\Dhat}{\widehat{D}}
\newcommand{\CC}{\mathrm{CC}}
\newcommand{\BP}{\mathrm{BP}}

\title{Uniform Mixing in Barrington Branching Programs}

\author{\vspace{3cm}\\Ernest Darell Plascencia Zerme\~no\\
\small Universidad Panamericana\\
\small \texttt{0244552@up.edu.mx}}
\date{}

\begin{document}

\maketitle

\vspace{2cm}
\begin{abstract}
We study the distribution $D_T$ over the symmetric group $S_5$ induced by evaluating
a Barrington branching program for the point function $\CC_T$ on uniformly random
inconsistent inputs. We prove two results. First, $D_T = D_{T'}$ for all
$T, T' \in \{0,1\}^\ell$ (exact T-independence), via a combinatorial flip bijection
that requires no properties of the underlying group. Second,
$\dtv(D_T, \mathrm{Uniform}(S_5)) \leq \negl(\ell)$ (spectral contraction),
via Fourier analysis on $S_5$ using the Diaconis--Shahshahani Upper Bound Lemma.
The contraction proof proceeds by exact computation of Fourier norms for the standard
representation $(4,1)$ at the first two levels of the commutator tree, followed by
an algebraic proof that the operator norm drops below~$1$ at the third level
--- enabled by the discovery that the level-$2$ obstruction has a regular simplex
structure in $\mathbb{R}^4$ that the commutator necessarily breaks.
The direction-mixing mechanism at level~$3$ is proved via the solvability of
$S_4 \cong \mathrm{Stab}(k)$: the derived series
$S_4 \rhd A_4 \rhd V_4 \rhd \{e\}$ forces the level-$3$ commutator
into the abelian Klein four-group, making it trivial within any single
stabilizer. We further prove that this spectral contraction mechanism is
\emph{specific to $S_5$}: for $S_n$ with $n \geq 6$, a dimensional
obstruction (the preserved subspace of $D_1$ has dimension $n-3$, and
two such subspaces in $\mathbb{R}^{n-1}$ must intersect when $n \geq 6$)
prevents operator-norm contraction at level~$3$. Together,
these results resolve the Uniform T-Independence Conjecture for Barrington
point-function branching programs, establish an information-theoretic foundation
for the security of branching-program-based obfuscation, and characterize
$S_5$ as the unique symmetric group where level-$3$ spectral contraction occurs.
\end{abstract}

\newpage
\tableofcontents
\newpage

%% ===================================================================
\section{Introduction}
\label{sec:intro}
%% ===================================================================

Indistinguishability obfuscation (iO) for general circuits is one of the most
powerful primitives in cryptography, implying a vast array of applications from
functional encryption to deniable encryption~\cite{GGH+13,SW14}. The first
candidate constructions~\cite{GGH+13} were based on matrix branching programs
over multilinear maps, randomized via Kilian's
technique~\cite{Kil88}, following the classical correspondence established by
Barrington's theorem~\cite{Bar89}: every $\mathrm{NC}^1$ circuit of depth $d$
can be simulated by a branching program of width~$5$ and length at most $4^d$
over the symmetric group $S_5$.

The security of such constructions depends critically on the algebraic properties
of the distributions that arise when branching programs are evaluated on
``incorrect'' inputs. Specifically, for a point function $\CC_T : \{0,1\}^\ell
\to \{0,1\}$ that outputs $1$ only on $x = T$, the Barrington program $\BP_T$
produces a designated permutation $\pi_{\mathrm{accept}} \in S_5$ on input $T$
and the identity $e$ on all other inputs. When the program is evaluated on a
random bit string $b \in \{0,1\}^L$ that is \emph{inconsistent} (at least one
variable receives different values at different steps), the output is a random
element of $S_5$ whose distribution we denote $D_T$.

Two properties of $D_T$ are essential for security:

\begin{enumerate}[label=(\roman*)]
\item \textbf{T-independence:} $D_T = D_{T'}$ for all $T, T'$, so the
  distribution reveals nothing about which target was chosen.
\item \textbf{Uniform mixing:} $\dtv(D_T, \mathrm{Uniform}(S_5)) \leq \negl(\ell)$,
  so the distribution is indistinguishable from random.
\end{enumerate}

These properties were widely assumed but not formally established:
the distributional behavior of Barrington programs on inconsistent inputs
--- a foundational question open since~$1989$ --- was implicitly relied upon
in security analyses without proof.
Property~(i) was used without justification, and property~(ii)
was supported by heuristic arguments and limited computational evidence.

\paragraph{Our contributions.}
We prove both properties.

For T-independence, we give a short combinatorial proof (\Cref{thm:t-indep})
via a flip bijection $F_S$ that maps inconsistent paths for $T$ to
inconsistent paths for $T'$ while preserving the output permutation. The proof
uses no property of $S_5$ and extends to conjunctions with wildcards.

For uniform mixing, we develop a Fourier-analytic framework on $S_5$
(\Cref{sec:contraction}). Using the Diaconis--Shahshahani Upper Bound
Lemma, we reduce the problem to bounding the Fourier norms of $D_T$ in
each irreducible representation. The standard representation $(4,1)$ of
dimension~$4$ is the bottleneck. We establish contraction at the first
two levels of the commutator tree by exact rational computation
(\Cref{thm:level01,thm:level12}), discover that the level-$2$ obstruction
has a regular simplex structure (\Cref{obs:simplex}), and prove that this
structure is necessarily broken at level~$3$ (\Cref{thm:level23}), yielding
operator-norm contraction with $\lambda \leq 0.0354$. An inductive argument
then gives super-exponential decay for all higher levels (\Cref{thm:induction}).

\paragraph{Implications.}
Combined, our results show that for $\ell \geq 8$, the distribution
over inconsistent evaluations of $\BP_T$ is simultaneously independent of
$T$ and negligibly close to uniform. This provides an information-theoretic
foundation for the security of obfuscation schemes based on Barrington branching
programs for point functions and, more broadly, for conjunctions with
bounded wildcard density.

\paragraph{Related work.}
The first iO candidate from branching programs~\cite{GGH+13} was followed
by constructions for specific function
classes: lockable obfuscation for point functions from
LWE~\cite{GKW17,WZ17}, conjunction obfuscation from noisy encodings
and LPN~\cite{BKMPRS18,BLMZ19}, and entropic Ring-LWE
approaches~\cite{BVWW16}.  The landmark result of Jain, Lin, and
Sahai~\cite{JLS21} achieved iO for general circuits from well-founded
assumptions, departing from the branching-program paradigm.  Recent
lattice-based candidates~\cite{BDGM20,WW21,HJL25} use matrix groups
$\mathrm{GL}_n(\mathbb{Z}_q)$ rather than $S_5$, while impossibility
results~\cite{DJMMV25} constrain the space of viable assumptions.
For these constructions, the mixed-input security model --- where an
adversary evaluates the program on inconsistent inputs --- remains
central~\cite{CVW18}.

Fourier-analytic methods for branching programs have a separate tradition
in the pseudorandomness literature~\cite{RSV13,LPV22}, but that
work uses Fourier analysis on $\{0,1\}^n$ (the input space) to fool
Boolean outputs of read-once programs.  Our analysis instead uses Fourier
analysis on $S_5$ (the output group) to bound total variation of the
group-valued output distribution of multi-read programs --- a fundamentally
different analytic object requiring different tools.  We are not aware of
prior work connecting the Diaconis--Shahshahani framework~\cite{DS81} to
branching-program security.

\paragraph{Techniques.}
The proof of T-independence is purely combinatorial. The proof of spectral
contraction uses three different methods at three different scales:
exact symbolic algebra (levels~$0$--$2$), an algebraic argument via
the solvability of $S_4$ (level~$3$), and analytic induction (levels~$4$+).
The key discovery that bridges levels~$2$ and~$3$ is the simplex structure of
the preserved directions: the $270$ out of $405$ level-$2$ configurations
with operator norm~$1$ each preserve one of exactly $5$ directions in
$\mathbb{R}^4$, forming a regular simplex with pairwise $\abs{\cos\theta} = 1/4$.
An algebraic proof shows that valid level-$3$ configurations cannot share a
preserved direction: within the stabilizer $\mathrm{Stab}(k) \cong S_4$,
the derived series terminates at depth~$3$, forcing the commutator to be
trivial. A separate analysis of $S_n$ for $n \geq 6$ reveals that this mechanism
is unique to $S_5$: the level-$1$ Fourier transform preserves a subspace of
dimension $n-3$ in $\mathbb{R}^{n-1}$, and the intersection bound
$2(n-3)-(n-1) = n-5 \geq 1$ blocks contraction for all $n \geq 6$.


%% ===================================================================
\section{Preliminaries}
\label{sec:prelim}
%% ===================================================================

\subsection{Branching Programs and Barrington's Theorem}

\begin{definition}[Permutation Branching Program]
A \emph{permutation branching program} of width $w$ and length $L$ over a group
$G \leq S_w$ is a sequence of instructions
$\Pi = ((i_1, g_1^{(0)}, g_1^{(1)}), \ldots, (i_L, g_L^{(0)}, g_L^{(1)}))$,
where each $i_j \in [\ell]$ is a variable index and $g_j^{(0)}, g_j^{(1)} \in G$.
On input $x \in \{0,1\}^\ell$, the program evaluates
$\Pi(x) = \prod_{j=1}^{L} g_j^{(x_{i_j})}$.
The program $\alpha$-\emph{computes} a function $f$ if $\Pi(x) = \alpha$ when
$f(x) = 1$ and $\Pi(x) = e$ when $f(x) = 0$.
\end{definition}

\begin{theorem}[Barrington~\cite{Bar89}]
\label{thm:barrington}
Every Boolean function computed by a circuit of depth $d$ can be
$\alpha$-computed by a permutation branching program of width $5$ and
length at most $4^d$ over $S_5$, for some $\alpha \neq e$.
\end{theorem}

The construction uses the \emph{commutator trick}: given programs $P$ and $Q$
that $\alpha$-compute $f$ and $\beta$-compute $g$ respectively, the concatenation
$P \cdot Q \cdot P^{-1} \cdot Q^{-1}$ computes $f \wedge g$ via the group
commutator $[\alpha, \beta] = \alpha\beta\alpha^{-1}\beta^{-1}$. This works because
$S_5$ is non-solvable: iterated commutators of non-identity elements remain
non-identity.

\begin{definition}[Point Function]
The point function $\CC_T : \{0,1\}^\ell \to \{0,1\}$ is defined by
$\CC_T(x) = 1$ if and only if $x = T$.
\end{definition}

For point functions, the Barrington program has $\ell$ leaves (one per bit position),
each comparing $x_i$ to $T_i$. Leaf $i$ outputs a fixed transposition $\sigma_i$
when $x_i = T_i$ and the identity $e$ when $x_i \neq T_i$. The tree of
commutators has depth $\log_2 \ell$, yielding a program of length $L = 4^{\log_2 \ell}
= \ell^2$.


\subsection{The Distribution $D_T$}

\begin{definition}[Inconsistent Path]
A bit string $b = (b_1, \ldots, b_L) \in \{0,1\}^L$ is \emph{inconsistent}
with respect to the program $\Pi$ if there exist steps $j, j'$ reading the
same variable ($i_j = i_{j'}$) with $b_j \neq b_{j'}$.
\end{definition}

\begin{definition}[$D_T$]
Let $\mathcal{I} \subset \{0,1\}^L$ be the set of inconsistent paths.
The distribution $D_T$ over $S_5$ is defined by
\[
D_T(\pi) = \Pr_{b \sim \mathrm{Uniform}(\mathcal{I})}
\left[ \prod_{j=1}^{L} g_j^{(b_j)} = \pi \right].
\]
\end{definition}


\subsection{Representation Theory of $S_5$}

The symmetric group $S_5$ has $7$ irreducible representations, labeled by
partitions of $5$:

\begin{center}
\begin{tabular}{lccl}
\toprule
Partition $\lambda$ & Dimension $d_\lambda$ & $\chi_\lambda(\text{transposition})$ & Name \\
\midrule
$(5)$       & 1 & 1 & Trivial \\
$(4,1)$     & 4 & 2 & Standard \\
$(3,2)$     & 5 & 1 & --- \\
$(3,1,1)$   & 6 & 0 & --- \\
$(2,2,1)$   & 5 & $-1$ & Sign $\otimes$ $(3,2)$ \\
$(2,1,1,1)$ & 4 & $-2$ & Sign $\otimes$ Standard \\
$(1^5)$     & 1 & $-1$ & Sign \\
\bottomrule
\end{tabular}
\end{center}

The dimensional check confirms: $1^2 + 4^2 + 5^2 + 6^2 + 5^2 + 4^2 + 1^2 = 120 = |S_5|$.

\begin{definition}[Standard Representation]
The standard representation $\rho_{(4,1)} : S_5 \to \mathrm{GL}(\mathbb{R}^4)$
acts on the hyperplane $V = \{v \in \mathbb{R}^5 : \sum_i v_i = 0\}$ by
permuting coordinates. We use the orthonormal basis
\begin{align*}
f_1 &= (e_0 - e_1)/\sqrt{2}, \quad
f_2 = (e_0 + e_1 - 2e_2)/\sqrt{6}, \\
f_3 &= (e_0 + e_1 + e_2 - 3e_3)/\sqrt{12}, \quad
f_4 = (e_0 + e_1 + e_2 + e_3 - 4e_4)/\sqrt{20},
\end{align*}
so that $\rho_{(4,1)}(\sigma)$ is orthogonal for all $\sigma \in S_5$.
\end{definition}


\subsection{The Diaconis--Shahshahani Upper Bound Lemma}

\begin{theorem}[Diaconis--Shahshahani~\cite{DS81}]
\label{thm:ds}
For any probability distribution $Q$ on a finite group $G$,
\[
4 \cdot \dtv(Q, U)^2 \leq \sum_{\rho \neq \mathrm{trivial}}
d_\rho \cdot \frob{\widehat{Q}(\rho)}^2,
\]
where the sum is over non-trivial irreducible representations $\rho$ of
dimension $d_\rho$, $\widehat{Q}(\rho) = \sum_{g \in G} Q(g) \rho(g)$ is the
Fourier transform, and $\frob{\cdot}$ denotes the Frobenius norm.
\end{theorem}

For $S_5$, the standard representation $(4,1)$ is the \emph{mixing
bottleneck}. Although $(2,1,1,1) = \mathrm{sign} \otimes (4,1)$ shares the
same normalized character ratio $\abs{\chi(\sigma)}/d_\rho = 1/2$, its leaf
projection $\frac{1}{2}(I - \rho_{(4,1)}(\sigma))$ has rank~$1$ (versus
rank~$3$ for $(4,1)$), causing it to contract at level~$1$ with operator
norm $\leq (1/2)^3 = 1/8$. The remaining representations satisfy
$\abs{\chi(\sigma)}/d_\rho \leq 1/5$ and likewise contract earlier
(\Cref{prop:other_reps}). It therefore suffices to
bound $\frob{\Dhat(4,1)}^2$.


%% ===================================================================
\section{Exact T-Independence}
\label{sec:t-indep}
%% ===================================================================

\begin{theorem}[T-Independence]
\label{thm:t-indep}
Let $\BP_T$ be the Barrington branching program for the point function $\CC_T$
over $\{0,1\}^\ell$. Then $D_T = D_{T'}$ for all $T, T' \in \{0,1\}^\ell$.
\end{theorem}

\begin{proof}
Let $S = \{i : T_i \neq T'_i\}$ be the set of differing bit positions.
Define $F_S : \{0,1\}^L \to \{0,1\}^L$ by
\[
F_S(b)_j =
\begin{cases}
1 - b_j & \text{if } i_j \in S, \\
b_j & \text{otherwise},
\end{cases}
\]
where $i_j$ is the variable read at step $j$.

We verify three properties:

\textbf{(1) $F_S$ is an involution.}
Applying $F_S$ twice flips each bit in $S$-positions twice, returning to
the original string. Hence $F_S \circ F_S = \Id$.

\textbf{(2) $F_S$ preserves inconsistency.}
$F_S$ flips \emph{all} readings of each variable $x_i$ with $i \in S$
simultaneously. If all readings of $x_i$ were equal (say, all $0$), after
flipping they are all $1$ --- still consistent. If some readings differed,
after flipping they still differ. Variables with $i \notin S$ are unaffected.
Therefore inconsistency of each variable is preserved individually.

\textbf{(3) $\varphi_{T'}(b) = \varphi_T(F_S(b))$ for all $b$.}
At each step $j$ reading variable $x_{i_j}$:
\begin{itemize}[nosep]
\item If $i_j \notin S$: $T_{i_j} = T'_{i_j}$ and $F_S$ does not flip $b_j$.
  Same input, same target bit, same output.
\item If $i_j \in S$: $T'_{i_j} = 1 - T_{i_j}$ and $F_S$ flips $b_j$.
  The swap of target bit and input bit cancel:
  $g_j^{(b_j)}$ under $T'$ equals $g_j^{(1-b_j)}$ under $T$, which equals
  $g_j^{(F_S(b)_j)}$ under $T$.
\end{itemize}
Hence the per-step outputs, and their product, are equal.

Combining (1)--(3):
\begin{align*}
D_{T'}(\pi) &= \Pr[\varphi_{T'}(b) = \pi \mid b \in \mathcal{I}] \\
&= \Pr[\varphi_T(F_S(b)) = \pi \mid b \in \mathcal{I}] \tag{by (3)} \\
&= \Pr[\varphi_T(b') = \pi \mid F_S^{-1}(b') \in \mathcal{I}]
\tag{substituting $b' = F_S(b)$} \\
&= \Pr[\varphi_T(b') = \pi \mid b' \in \mathcal{I}]
\tag{by (2), since $F_S = F_S^{-1}$ by (1)} \\
&= D_T(\pi). \qedhere
\end{align*}
\end{proof}

\begin{remark}
The proof uses no property of $S_5$ --- it works for branching programs of
any width over any group. It applies to all programs for point functions,
not just Barrington's construction.
\end{remark}

\begin{remark}
The proof extends to conjunctions with wildcards: if position $i$ is a wildcard,
the corresponding leaf outputs the same permutation $\sigma_i$ for both bit values.
Flipping $b_j$ at wildcard positions changes both branches identically, so the output
is unchanged. The bijection $F_S$ is restricted to non-wildcard differing positions.
\end{remark}

\begin{remark}[Unconditional T-independence]
\label{rem:unconditional}
The bijection $F_S$ preserves inconsistency of each variable individually
(property~(2)), and therefore also preserves the set of \emph{consistent}
paths. Consequently, $D_T = D_{T'}$ holds not only for the distribution
conditioned on inconsistent paths, but also for the unconditional distribution
over all $2^L$ paths and for the distribution over consistent paths alone.
The conditioning on inconsistency is needed for the security argument
(consistent paths reveal~$T$), not for T-independence itself.
\end{remark}

\begin{remark}[Variable-grouping invariance]
\label{rem:grouping}
The distribution $D_T$ depends on the choice of transpositions assigned to
each leaf, but \emph{not} on how variables are grouped into the binary
commutator tree. Exhaustive computation for $\ell = 4$ confirms that all
three possible balanced groupings of $4$ variables yield identical distributions.
This is because $F_S$ acts on variables, not on tree positions.
\end{remark}


%% ===================================================================
\section{Spectral Contraction}
\label{sec:contraction}
%% ===================================================================

By \Cref{thm:t-indep}, $D_T$ is independent of $T$, so it suffices to show
$\dtv(D_T, U) \leq \negl(\ell)$ for any fixed $T$. By the
Diaconis--Shahshahani bound (\Cref{thm:ds}), this reduces to bounding
$\frob{\Dhat(\rho)}^2$ for each non-trivial representation $\rho$.

\begin{remark}[Sign representation vanishes identically]
\label{rem:sign}
For the sign representation $(1^5)$, the leaf Fourier transform is
$\Dhat(\mathrm{sign}) = \frac{1}{2}(\mathrm{sign}(\sigma) + 1) = 0$
for any transposition~$\sigma$, since $\mathrm{sign}(\sigma) = -1$.
This propagates to all levels by multiplicativity, so the sign
representation contributes nothing to the Diaconis--Shahshahani bound.
The sum thus runs over $5$ (not $6$) non-trivial representations.
\end{remark}

\begin{proposition}[Contraction in the remaining representations]
\label{prop:other_reps}
For all five non-trivial, non-sign representations $\rho$ of $S_5$, the
operator norm $\opn{\Dhat_3(\rho)}$ is strictly less than $1$ for every
valid level-$3$ configuration. The representation $(4,1)$ is the unique
bottleneck:
\begin{center}
\begin{tabular}{lcccl}
\toprule
$\rho$ & $d_\rho$ & Leaf rank & $\max\opn{\Dhat_3(\rho)}$ & First level $< 1$ \\
\midrule
$(4,1)$     & 4 & 3 & $3.54 \times 10^{-2}$ & Level $3$ \\
$(3,2)$     & 5 & 3 & $1.77 \times 10^{-5}$ & Level $2$ \\
$(3,1,1)$   & 6 & 3 & $1.05 \times 10^{-5}$ & Level $2$ \\
$(2,2,1)$   & 5 & 2 & $< 10^{-10}$          & Level $1$ \\
$(2,1,1,1)$ & 4 & 1 & $< 10^{-10}$          & Level $1$ \\
\bottomrule
\end{tabular}
\end{center}
The representations $(2,1,1,1) = \mathrm{sign} \otimes (4,1)$ and
$(2,2,1) = \mathrm{sign} \otimes (3,2)$ have leaf projections of rank $1$
and $2$ respectively. Their level-$1$ operator norms are at most $(1/2)^3 = 1/8$,
so they contract immediately. The representations $(3,2)$ and $(3,1,1)$
have rank-$3$ leaf projections like $(4,1)$, but lack the simplex
obstruction (\Cref{obs:simplex}): all $405$ level-$2$ quadruples have
operator norm $< 0.085$, so they contract one level earlier.
\emph{Verified exhaustively over all $73810$ level-$3$ configurations
for each representation} (Script~$13$ in the accompanying code).
\end{proposition}

We focus on the bottleneck representation $(4,1)$.


\subsection{Framework: Projections and the Fourier Domain}
\label{sec:framework}

For a transposition $\sigma = (a\;b) \in S_5$, the standard representation
$\rho_{(4,1)}(\sigma)$ has eigenvalues $\{+1, +1, +1, -1\}$ (three fixed
directions and one flipped). The Fourier transform of the leaf distribution
$\{\sigma : 1/2,\; e : 1/2\}$ is
\[
P_\sigma \coloneqq \frac{1}{2}\bigl(\rho_{(4,1)}(\sigma) + I_4\bigr),
\]
which is a rank-$3$ orthogonal projection (idempotent, symmetric, trace $3$).

\begin{lemma}
\label{lem:proj_basic}
For any transposition $\sigma \in S_5$:
\begin{enumerate}[label=(\alph*),nosep]
\item $P_\sigma^2 = P_\sigma$ (idempotent),
\item $P_\sigma^\top = P_\sigma$ (symmetric),
\item $\mathrm{rank}(P_\sigma) = 3$,
\item $\frob{P_\sigma}^2 = \tr(P_\sigma) = 3$.
\end{enumerate}
\end{lemma}

The null space of $P_\sigma$ is spanned by the projection of $e_a - e_b$
onto the hyperplane $V$, where $\sigma = (a\;b)$.

At each level of the Barrington commutator tree, the Fourier transform is a product
of projections. For the single AND block at level~$1$ using transpositions $\sigma, \tau$
(with $\sigma^{-1} = \sigma$, $\tau^{-1} = \tau$ since they are transpositions):
\[
\Dhat_1 = P_\sigma \cdot P_\tau \cdot P_\sigma \cdot P_\tau = (P_\sigma P_\tau)^2.
\]
At level~$2$, using two level-$1$ blocks $A$ (with $\sigma_1, \tau_1$) and
$B$ (with $\sigma_2, \tau_2$):
\[
\Dhat_2 = (P_{\sigma_1} P_{\tau_1})^2 \cdot (P_{\sigma_2} P_{\tau_2})^2
\cdot (P_{\tau_1} P_{\sigma_1})^2 \cdot (P_{\tau_2} P_{\sigma_2})^2,
\]
where the inverse blocks reverse the transposition order.


\subsection{Theorem 1: Level 0 to Level 1 Contraction}

\begin{theorem}[Universal Level-$1$ Contraction]
\label{thm:level01}
For any pair of adjacent transpositions $(\sigma, \tau)$ in $S_5$
(i.e., $[\sigma, \tau] \neq e$):
\[
\frob{(P_\sigma P_\tau)^2}^2 = \frac{129}{64}.
\]
The contraction ratio is $\frac{129/64}{3} = \frac{43}{64} \approx 0.6719$,
and is \textbf{universal} --- identical for all $30$ adjacent transposition pairs.
\end{theorem}

\begin{proof}
There are $\binom{5}{2} = 10$ transpositions in $S_5$, forming $30$ adjacent
pairs (pairs sharing exactly one element, equivalently $[\sigma,\tau] \neq e$).
For each pair, we compute $(P_\sigma P_\tau)^2$ in the orthonormal basis of
\Cref{sec:prelim} using exact rational arithmetic, and verify
$\frob{(P_\sigma P_\tau)^2}^2 = 129/64$.

The universality follows from the geometry of the standard representation:
the null vectors of $P_\sigma$ and $P_\tau$ (projections of $e_a - e_b$ and
$e_c - e_d$ onto $V$) have inner product $\pm 1/2$ for all adjacent pairs,
determined by the simplex geometry of $\{e_0, \ldots, e_4\}$ in $V$.
Since $\frob{(P_\sigma P_\tau)^2}^2$ depends only on this angle, it is
universal.
\end{proof}


\subsection{Theorem 2: Level 1 to Level 2 Contraction}

\begin{theorem}[Level-$2$ Contraction]
\label{thm:level12}
For any valid quadruple of transpositions $(\sigma_1, \tau_1, \sigma_2, \tau_2)$
with $[\sigma_i, \tau_i] \neq e$ and $[[\sigma_1,\tau_1], [\sigma_2,\tau_2]] \neq e$:
\[
\frob{\Dhat_2}^2 \leq \frac{67586345}{67108864} \approx 1.0071.
\]
The contraction ratio relative to level $1$ is at most
$\frac{67586345/67108864}{129/64} \approx 0.4997 < 1$.
\end{theorem}

\begin{proof}
We enumerate all $405$ valid quadruples (from the $30$ adjacent pairs with the
additional constraint that the outer commutator is nontrivial). For each quadruple,
we compute $\Dhat_2$ in exact rational arithmetic and evaluate $\frob{\Dhat_2}^2$.
The worst case over all $405$ quadruples is $67586345/67108864$.
The computation yields $16$ distinct Frobenius norms, ranging from
$82431849/1073741824 \approx 0.077$ to the worst case above.
\end{proof}


\subsection{The Simplex Obstruction}
\label{sec:simplex}

Despite the Frobenius-norm contraction at level~$2$, the operator norm
$\opn{\Dhat_2}$ equals $1$ for $270$ of the $405$ valid quadruples,
blocking a naive submultiplicativity argument for higher levels.
The following observation explains the structure of this obstruction.

\begin{observation}[Simplex Structure]
\label{obs:simplex}
The $270$ quadruples with $\opn{\Dhat_2} = 1$ each preserve exactly one direction
in $\mathbb{R}^4$ (as the eigenvector of $\Dhat_2^\top \Dhat_2$ with eigenvalue $1$).
These directions are precisely the projections of the $5$ standard basis vectors
$e_0, \ldots, e_4 \in \mathbb{R}^5$ onto the hyperplane $V = \{\sum v_i = 0\}$:
\[
v_k = \frac{e_k - \bar{e}}{\norm{e_k - \bar{e}}}, \quad
\bar{e} = \tfrac{1}{5}(1,1,1,1,1)^\top, \quad k = 0, \ldots, 4.
\]
These $5$ directions form a \textbf{regular simplex} in $\mathbb{R}^4$, with
pairwise inner products $\abs{\langle v_i, v_j \rangle} = 1/4$ for all $i \neq j$.
Each direction is preserved by exactly $54$ quadruples (perfectly symmetric).
\end{observation}

The regularity of this structure --- five directions, uniform distribution,
constant pairwise angle --- suggests a general phenomenon rather than
an artifact of~$S_5$.

\begin{theorem}[Dimensional obstruction for $S_n$, $n \geq 6$]
\label{thm:dim-obstruction}
For $S_n$ with $n \geq 5$, the level-$1$ Fourier transform
$D_1 = (P_\sigma P_\tau)^2$ in the standard representation $(n-1,1)$
has singular values $\underbrace{1, \ldots, 1}_{n-3},\; 1/8,\; 0$.
In particular, $D_1$ preserves a subspace of dimension $n-3$ isometrically.

For $n \geq 6$, the minimum intersection of two such preserved subspaces
in $\mathbb{R}^{n-1}$ is $2(n-3)-(n-1) = n-5 \geq 1$, so every level-$2$
quadruple has $\opn{\Dhat_2} = 1$ and level-$3$ contraction fails.
Consequently, the spectral contraction mechanism of this paper is specific
to $S_5$.
\end{theorem}

\begin{proof}
The projections $P_\sigma = I_V - v_\sigma v_\sigma^\top$ and
$P_\tau = I_V - v_\tau v_\tau^\top$ share the property that
$\langle v_\sigma, v_\tau \rangle = -1/2$ for all adjacent pairs
(independent of~$n$). Decomposing $V = W \oplus W^\perp$ where
$W = \mathrm{span}(v_\sigma, v_\tau)$ (dimension~$2$):
on $W^\perp$ (dimension $n-3$), both projections act as the identity,
giving $n-3$ singular values equal to~$1$;
on~$W$, the $2\times 2$ block $(P_\sigma P_\tau)^2\big|_W$ has
$D_1^\top D_1\big|_W$ with eigenvalues $1/64$ and~$0$
(trace $= 4/256$, determinant $= 0$), giving singular values $1/8$ and~$0$.
The $2\times 2$ block is independent of~$n$.

For $n \geq 6$: two subspaces of dimension $n-3$ in $\mathbb{R}^{n-1}$
intersect in dimension $\geq n-5 \geq 1$. Any unit vector in this
intersection is preserved isometrically by both level-$1$ blocks, so
$\opn{\Dhat_2} = 1$ for every valid quadruple.
\end{proof}

\begin{remark}
The simplex structure (\Cref{obs:simplex}) does not generalize to~$S_n$
for $n \geq 6$: exhaustive computation confirms that all $1620$ valid
quadruples for $S_6$ and all $4725$ for $S_7$ have operator norm~$1$
(versus $270/405 = 66.7\%$ for $S_5$), and the preserved directions
do not form a regular simplex. The dimensional obstruction explains this:
$S_5$ is the unique case where the preserved subspace dimension ($n-3 = 2$)
allows disjoint placement in the ambient space ($n-1 = 4$).
\end{remark}


\subsection{Theorem 3: Level 2 to Level 3 Contraction}

\begin{lemma}[Direction Mixing]
\label{lem:mixing}
For every valid level-$3$ configuration $(\mathrm{quad}_A, \mathrm{quad}_B)$
in which both sub-blocks have $\opn{\Dhat_2} = 1$, the two sub-blocks
preserve \textbf{different} directions.
\end{lemma}

\begin{proof}
We establish two claims.

\textbf{Claim 1:} A level-$2$ quadruple $(\sigma_1, \tau_1, \sigma_2, \tau_2)$
with $\opn{\Dhat_2} = 1$ preserves direction $v_k$ if and only if all four
transpositions fix element~$k$.

The forward direction is immediate: if all transpositions fix~$k$, then
$\langle v_k, v_{(a,b)} \rangle = (\delta_{k,a} - \delta_{k,b})/(\sqrt{2}\,\norm{e_k - \bar{e}}) = 0$
for each null vector $v_{(a,b)}$, so $v_k \in W_A^\perp \cap W_B^\perp$ and $\Dhat_2 v_k = v_k$.
For the backward direction, if $\opn{\Dhat_2 v_k} = 1$ and $\Dhat_2 = D_{1,A} D_{1,B} D_{1,A}^\top D_{1,B}^\top$
with each factor having operator norm $\leq 1$, the chain of norm inequalities forces
$v_k$ into the preserved subspace $W^\perp$ of each level-$1$ block, hence
$\langle v_k, v_{(a,b)} \rangle = 0$ for all four transpositions, i.e., all fix~$k$.

\textbf{Claim 2:} If both quadruples have all transpositions fixing element~$k$,
the level-$3$ commutator is trivial.

All eight transpositions lie in $\mathrm{Stab}(k) \cong S_4$, which is
\emph{solvable} with derived series
$S_4 \rhd A_4 \rhd V_4 \rhd \{e\}$,
where $V_4 \cong \mathbb{Z}_2 \times \mathbb{Z}_2$ is the Klein four-group.
Following the Barrington commutator tree within $\mathrm{Stab}(k)$:
level-$1$ targets (commutators of adjacent transpositions) are $3$-cycles
in $[S_4, S_4] = A_4$;
level-$2$ targets (commutators of $3$-cycles) lie in $[A_4, A_4] = V_4$;
the level-$3$ commutator $[\alpha_A, \alpha_B]$ with $\alpha_A, \alpha_B \in V_4$
lies in $[V_4, V_4] = \{e\}$, since $V_4$ is abelian.

This contradicts the requirement $[\alpha_A, \alpha_B] \neq e$ for a valid
level-$3$ configuration. Therefore both quadruples cannot preserve the same
direction.
\end{proof}

\begin{theorem}[Level-$3$ Contraction]
\label{thm:level23}
For every valid level-$3$ configuration:
\[
\opn{\Dhat_3} \leq \lambda \coloneqq 0.0354.
\]
In particular, $\opn{\Dhat_3} < 1$ for all $73810$ configurations.
\end{theorem}

\begin{proof}
For each valid level-$3$ pair $(\mathrm{quad}_A, \mathrm{quad}_B)$, we compute
\[
\Dhat_3 = \Dhat_{2,A} \cdot \Dhat_{2,B} \cdot \Dhat_{2,A}^\top \cdot \Dhat_{2,B}^\top
\]
and evaluate the operator norm (maximum singular value) numerically. The maximum
over all $73810$ configurations is $\lambda \approx 0.0353$, strictly less than $1$.

The contraction occurs because:
\begin{itemize}[nosep]
\item If at least one sub-block has $\opn{\Dhat_2} < 1$, submultiplicativity gives
  $\opn{\Dhat_3} \leq \opn{\Dhat_{2,A}}^2 \cdot \opn{\Dhat_{2,B}}^2 < 1$.
\item If both sub-blocks have $\opn{\Dhat_2} = 1$ but preserve different directions
  $v_A \neq v_B$ (guaranteed by \Cref{lem:mixing}), then $\Dhat_{2,A}$ maps $v_B$
  to a vector with component $\abs{\langle v_A, v_B \rangle} = 1/4$ along $v_A$,
  causing the product to contract. \qedhere
\end{itemize}
\end{proof}


\subsection{Theorem 4: Induction to All Levels}

\begin{theorem}[Exponential Contraction]
\label{thm:induction}
For all Barrington branching programs for point functions $\CC_T$ over
$\{0,1\}^\ell$ with $\ell = 2^d$, $d \geq 3$:
\[
\opn{\Dhat_d\bigl((4,1)\bigr)} \leq \lambda^{4^{d-3}},
\]
where $\lambda = 0.0354$.
\end{theorem}

\begin{proof}
At depth $d+1$, the commutator structure gives
$\Dhat_{d+1} = \Dhat_A \cdot \Dhat_B \cdot \Dhat_A^\top \cdot \Dhat_B^\top$.
By submultiplicativity of the operator norm and $\opn{M^\top} = \opn{M}$:
\[
\opn{\Dhat_{d+1}} \leq \opn{\Dhat_A} \cdot \opn{\Dhat_B} \cdot
\opn{\Dhat_A^\top} \cdot \opn{\Dhat_B^\top} = \opn{\Dhat_d}^4.
\]
The base case $\opn{\Dhat_3} \leq \lambda$ is \Cref{thm:level23}.
Iterating: $\opn{\Dhat_d} \leq \lambda^{4^{d-3}}$.
\end{proof}


%% ===================================================================
\section{Main Result}
\label{sec:main}
%% ===================================================================

\begin{theorem}[Main Theorem]
\label{thm:main}
For all Barrington branching programs for point functions $\CC_T$ over
$\{0,1\}^\ell$ with $\ell = 2^d$, $d \geq 3$:
\begin{enumerate}[label=(\alph*),nosep]
\item $D_T = D_{T'}$ for all $T, T' \in \{0,1\}^\ell$.
\item $\dtv(D_T, \mathrm{Uniform}(S_5)) \leq 6 \cdot \lambda^{4^{d-3}}$,
  where $\lambda = 0.0354$.
\item In particular, $\dtv(D_T, U) = \negl(\ell)$.
\end{enumerate}
\end{theorem}

\begin{proof}
Part (a) is \Cref{thm:t-indep}. For part (b), by \Cref{thm:ds}:
\[
\dtv(D_T, U)^2 \leq \frac{1}{4} \sum_{\rho \neq \mathrm{triv}} d_\rho \cdot
\frob{\Dhat(\rho)}^2.
\]
By \Cref{rem:sign}, the sign representation contributes zero.
For each remaining representation $\rho$ of dimension $d_\rho$,
\Cref{prop:other_reps} gives a base case
$\opn{\Dhat_3(\rho)} \leq \lambda_\rho < 1$ for all valid level-$3$
configurations. The submultiplicativity argument of \Cref{thm:induction}
applies to each $\rho$ independently (it uses only $\opn{M^\top} = \opn{M}$
and submultiplicativity, which hold in any representation), yielding
$\opn{\Dhat_d(\rho)} \leq \lambda_\rho^{4^{d-3}}$.
Since $\frob{M}^2 \leq d_\rho \cdot \opn{M}^2$:
\[
\sum_{\rho} d_\rho \cdot \frob{\Dhat(\rho)}^2
\leq \sum_{\rho} d_\rho^2 \cdot \lambda_\rho^{2 \cdot 4^{d-3}}
\leq \Bigl(\sum_{\rho} d_\rho^2\Bigr) \cdot \lambda^{2 \cdot 4^{d-3}},
\]
where $\lambda = \max_\rho \lambda_\rho = \lambda_{(4,1)} \approx 0.0353$
is the bottleneck (\Cref{prop:other_reps}). Since
$\sum_\rho d_\rho^2 = 4^2 + 5^2 + 6^2 + 5^2 + 4^2 = 118$:
\[
\dtv(D_T, U)^2 \leq \frac{118}{4} \cdot \lambda^{2 \cdot 4^{d-3}},
\qquad
\dtv(D_T, U) \leq \sqrt{\tfrac{118}{4}} \cdot \lambda^{4^{d-3}}
< 6 \cdot \lambda^{4^{d-3}}.
\]
For part (c), since $\lambda < 1$ and $4^{d-3}$ grows super-exponentially in
$d = \log_2 \ell$, the bound is negligible in $\ell$.
\end{proof}

For concrete values:
\begin{center}
\begin{tabular}{cccc}
\toprule
$\ell$ & Depth $d$ & $\opn{\Dhat_{(4,1)}}$ bound & $\dtv$ bound \\
\midrule
$8$   & $3$ & $3.54 \times 10^{-2}$ & $2.12 \times 10^{-1}$ \\
$16$  & $4$ & $1.57 \times 10^{-6}$ & $9.42 \times 10^{-6}$ \\
$32$  & $5$ & $6.08 \times 10^{-24}$ & $3.65 \times 10^{-23}$ \\
$64$  & $6$ & $1.37 \times 10^{-93}$ & $8.21 \times 10^{-93}$ \\
\bottomrule
\end{tabular}
\end{center}


%% ===================================================================
\section{Extensions and Open Problems}
\label{sec:extensions}
%% ===================================================================

\subsection{Conjunctions with Wildcards}

A conjunction $\CC_{T,W}$ with target $T \in \{0,1\}^\ell$ and wildcard set
$W \subseteq [\ell]$ outputs $1$ iff $x_i = T_i$ for all $i \notin W$.
In the branching program, wildcard positions output $\sigma_i$ for both
bit values.  Obfuscation of conjunctions has been studied
extensively~\cite{BVWW16,BKMPRS18,BLMZ19}; our results provide an
information-theoretic complement to the computational security guarantees
in that literature.

\Cref{thm:t-indep} extends directly to conjunctions:
$D_{T,W} = D_{T',W}$ for all $T, T'$ with the same wildcard set $W$.
The spectral contraction depends on the wildcard density $|W|/\ell$.
Computational evidence shows strong contraction for $|W|/\ell \leq 1/2$,
weakening contraction for $|W|/\ell = 3/4$, and no contraction for $|W|/\ell = 1$
(the trivial function). The precise threshold is an open problem.


\subsection{Beyond Point Functions}

For general branching programs computing the same Boolean function via
different constructions, ``strong universal mixing'' ($D_{C_0} = D_{C_1}$)
is \textbf{false}: programs with different transposition assignments produce
different distributions. However, ``weak universal mixing'' (both distributions
converge to uniform independently) appears to hold: computational evidence shows
equivalent circuits with different pair assignments achieving similar $\dtv$
values. Proving this would strengthen the security guarantee for
indistinguishability obfuscation of general $\mathrm{NC}^1$ functions.

Computational experiments reveal a stronger pattern.  For $\ell = 4$, two
programs computing $\mathrm{AND}_4$ with different transposition pairs
yield Fourier norms $\frob{\Dhat(4,1)}^2 = 0.147$ and $0.293$
respectively --- clearly different distributions --- yet \emph{identical}
total variation distances $\dtv = 0.309$ to uniform.  The same equidistance
persists for OR-of-AND circuits ($\dtv \approx 0.19$ for two structurally
different implementations).  If this ``equidistance'' phenomenon is general,
it would mean that $\dtv(D_C, U)$ depends only on the Boolean function
computed, not on the specific Barrington construction --- a property
strictly stronger than weak universal mixing.


\subsection{Generalization to Other Groups}

\Cref{thm:dim-obstruction} shows that our level-$3$ contraction mechanism
does not extend to $S_n$ for $n \geq 6$. The dimensional obstruction is
fundamental: in the standard representation of $S_n$, the level-$1$
Fourier transform preserves $n-3$ of $n-1$ dimensions, and the ratio
$(n-3)/(n-1) \to 1$ as $n \to \infty$, leaving insufficient room for
contraction.

This does not rule out uniform mixing for $S_n$ with $n > 5$ via
alternative proof strategies. Two natural directions remain open:
\begin{enumerate}[nosep]
\item \emph{Higher-level contraction.} Although level-$3$ contraction
  fails for $n \geq 6$, it is conceivable that contraction occurs at
  level~$4$ or higher, where the commutator structure may break the
  shared preserved directions. Computational verification of this
  hypothesis for $S_6$ is feasible but requires enumerating
  $O(10^{12})$ level-$4$ configurations.
\item \emph{Non-standard representations.} For $S_n$ with $n > 5$,
  non-standard representations (e.g., $(n-2,2)$ or $(n-2,1,1)$) may
  exhibit better contraction properties. The bottleneck analysis
  (\Cref{prop:other_reps}) would need to be repeated for each~$n$.
\end{enumerate}

For the matrix groups $\mathrm{GL}_n(\mathbb{Z}_q)$ used in modern
lattice-based iO constructions~\cite{BDGM20,WW21,HJL25}, the
representation theory is richer (principal series, cuspidal, Steinberg
representations), and the dimensional obstruction may not apply in the
same form. Determining whether spectral contraction is achievable for
$\mathrm{GL}_2(\mathbb{F}_q)$ with small~$q$ remains a concrete open problem,
especially timely given recent impossibility results for evasive
LWE~\cite{DJMMV25}.


\subsection{Why $S_5$ is Structurally Unique}

The results of this paper reveal a precise structural explanation for
why $S_5$ is the unique symmetric group where level-$3$ spectral
contraction succeeds. Three properties conspire:

\begin{enumerate}[nosep]
\item \textbf{$S_5$ is non-solvable.} $A_5$ is simple, so iterated
  commutators of non-identity elements never reach the identity ---
  this is what makes Barrington's theorem work over~$S_5$.
\item \textbf{$\mathrm{Stab}(k) \cong S_4$ is solvable with derived
  length~$3$.} The derived series
  $S_4 \rhd A_4 \rhd V_4 \rhd \{e\}$ terminates at depth~$3$,
  which is precisely the level of the commutator tree where contraction
  occurs. This forces the direction-mixing mechanism (\Cref{lem:mixing}).
\item \textbf{The preserved subspace dimension equals half the ambient
  dimension.} $D_1$ preserves $n-3 = 2$ of $n-1 = 4$ dimensions,
  allowing two preserved subspaces to be placed in general position
  (intersection dimension $= 0$). For $n \geq 6$, the ratio
  $(n-3)/(n-1) > 1/2$ makes this impossible (\Cref{thm:dim-obstruction}).
\end{enumerate}

The coincidence of these three properties --- non-solvability of $S_5$,
solvability of $S_4$ at derived length exactly~$3$, and the critical
dimension ratio --- makes $S_5$ the unique locus where this proof strategy
succeeds.


%% ===================================================================
\section{Computational Verification}
\label{sec:computation}
%% ===================================================================

The proofs of \Cref{thm:level01,thm:level12} use exact rational arithmetic
(implemented in \texttt{sympy}) with a verified orthonormal basis for the
standard representation. \Cref{thm:level23} uses numerical computation
(\texttt{numpy}) over all $73810$ valid configurations.

The full verification suite consists of $15$ scripts,
publicly available at:
\begin{center}
\url{https://github.com/Jefemaestro33/barrington-mixing}
\end{center}

\begin{center}
\begin{tabular}{clc}
\toprule
Script & Content & Method \\
\midrule
01 & Ambivalence of $S_5$ (similarity blindness) & Exhaustive \\
02 & Collision analysis, $\ell = 4$ & Exhaustive \\
03 & Collision scaling, $\ell = 4 \to 8$ & Sampled \\
04 & Fourier spectral analysis & Exact + sampled \\
05 & Barrington construction + Kilian obfuscation & Exhaustive \\
06 & Contraction operator (unconditional distribution) & Analytic \\
07 & T-independence proof verification & Exhaustive \\
08 & Conjunctions with wildcards & Exhaustive \\
09 & Universal mixing investigation & Exact + sampled \\
10 & Complete spectral contraction proof & Exact + numerical \\
11 & Eigenvector analysis (simplex discovery) & Numerical \\
12 & Level-$3$ exhaustive verification & Numerical \\
13 & All-representations contraction verification & Numerical \\
14 & $S_6$/$S_7$ dimensional obstruction (\Cref{thm:dim-obstruction}) & Exact + exhaustive \\
15 & Algebraic direction mixing verification (\Cref{lem:mixing}) & Exact \\
\bottomrule
\end{tabular}
\end{center}

Scripts 10--13 collectively verify the main theorem in under $20$ seconds
on a standard laptop. Script~13 verifies \Cref{prop:other_reps}:
operator-norm contraction for all five non-trivial, non-sign representations
across all $73810$ level-$3$ configurations.
Script~14 verifies \Cref{thm:dim-obstruction}: the dimensional obstruction
for $S_6$ ($1620$ quadruples, all with operator norm~$1$) and $S_7$
($4725$ quadruples), confirming that level-$3$ contraction fails for $n \geq 6$.
Script~15 verifies each step of the algebraic proof of \Cref{lem:mixing}:
the equivalence between preserved directions and fixed elements, the derived
series of $S_4$, and the triviality of $[V_4, V_4]$.
Scripts~01--09 provide additional cross-validation and extensions.


%% ===================================================================
%% References
%% ===================================================================

\clearpage

\begin{thebibliography}{99}

\bibitem{Bar89}
D.~A.~Barrington.
\newblock Bounded-width polynomial-size branching programs recognize exactly
  those languages in $\mathrm{NC}^1$.
\newblock \emph{J.~Comput.~Syst.~Sci.}, 38(1):150--164, 1989.

\bibitem{AGIS14}
P.~Ananth, D.~Gupta, Y.~Ishai, and A.~Sahai.
\newblock Optimizing obfuscation: Avoiding Barrington's theorem.
\newblock In \emph{CCS}, pages 646--658, 2014.

\bibitem{BDGM20}
Z.~Brakerski, N.~D{\"o}ttling, S.~Garg, and G.~Malavolta.
\newblock Candidate iO from homomorphic encryption schemes.
\newblock In \emph{EUROCRYPT}, pages 79--109, 2020.

\bibitem{BVWW16}
Z.~Brakerski, V.~Vaikuntanathan, H.~Wee, and D.~Wichs.
\newblock Obfuscating conjunctions under entropic ring LWE.
\newblock In \emph{ITCS}, 2016.

\bibitem{BKMPRS18}
A.~Bishop, L.~Kowalczyk, T.~Malkin, V.~Pastro, M.~Raykova, and K.~Shi.
\newblock A simple obfuscation scheme for pattern-matching with wildcards.
\newblock In \emph{CRYPTO}, pages 731--752, 2018.

\bibitem{BLMZ19}
J.~Bartusek, T.~Lepoint, F.~Ma, and M.~Zhandry.
\newblock New techniques for obfuscating conjunctions.
\newblock In \emph{EUROCRYPT}, pages 636--666, 2019.

\bibitem{DS81}
P.~Diaconis and M.~Shahshahani.
\newblock Generating a random permutation with random transpositions.
\newblock \emph{Z.~Wahrscheinlichkeitstheorie verw.~Gebiete}, 57:159--179, 1981.

\bibitem{CVW18}
Y.~Chen, V.~Vaikuntanathan, and H.~Wee.
\newblock GGH15 beyond permutation branching programs: Proofs, attacks, and
  candidates.
\newblock In \emph{CRYPTO}, pages 577--607, 2018.

\bibitem{DJMMV25}
N.~D{\"o}ttling, A.~Jain, G.~Malavolta, S.~Mathialagan, and
  V.~Vaikuntanathan.
\newblock Simple and general counterexamples for private-coin evasive LWE.
\newblock In \emph{CRYPTO}, 2025.

\bibitem{GGH+13}
S.~Garg, C.~Gentry, S.~Halevi, M.~Raykova, A.~Sahai, and B.~Waters.
\newblock Candidate indistinguishability obfuscation and functional encryption
  for all circuits.
\newblock In \emph{FOCS}, pages 40--49, 2013.

\bibitem{GKW17}
R.~Goyal, V.~Koppula, and B.~Waters.
\newblock Lockable obfuscation.
\newblock In \emph{FOCS}, pages 612--621, 2017.

\bibitem{WZ17}
D.~Wichs and G.~Zirdelis.
\newblock Obfuscating compute-and-compare programs under LWE.
\newblock In \emph{FOCS}, pages 600--611, 2017.

\bibitem{HJL25}
Y.-C.~Hsieh, A.~Jain, and H.~Lin.
\newblock Lattice-based post-quantum iO from circular security with random
  opening assumption.
\newblock In \emph{CRYPTO}, 2025.

\bibitem{JLS21}
A.~Jain, H.~Lin, and A.~Sahai.
\newblock Indistinguishability obfuscation from well-founded assumptions.
\newblock In \emph{STOC}, pages 60--73, 2021.

\bibitem{Kil88}
J.~Kilian.
\newblock Founding cryptography on oblivious transfer.
\newblock In \emph{STOC}, pages 20--31, 1988.

\bibitem{LPV22}
C.~H.~Lee, E.~Pyne, and S.~Vadhan.
\newblock Fourier growth of regular branching programs.
\newblock In \emph{RANDOM}, 2022.

\bibitem{RSV13}
O.~Reingold, T.~Steinke, and S.~Vadhan.
\newblock Pseudorandomness for regular branching programs via Fourier analysis.
\newblock In \emph{RANDOM}, pages 655--670, 2013.

\bibitem{SW14}
A.~Sahai and B.~Waters.
\newblock How to use indistinguishability obfuscation: deniable encryption, and
  more.
\newblock In \emph{STOC}, pages 475--484, 2014.

\bibitem{WW21}
H.~Wee and D.~Wichs.
\newblock Candidate obfuscation via oblivious LWE sampling.
\newblock In \emph{EUROCRYPT}, 2021.

\end{thebibliography}

\end{document}